% =====================================================================
% RSL: Física computacional y gestión del riesgo de desastres
% Curso: Organización y Métodos (UNSA, 2026-B) -- Mg. William Bornás Rios
% Compilar:  pdflatex -> biber -> pdflatex -> pdflatex   (o: latexmk -pdf main)
% Los marcadores \relleno{...} indican texto pendiente.
% =====================================================================
\documentclass[12pt,a4paper]{article}

\usepackage[utf8]{inputenc}
\usepackage[T1]{fontenc}
\usepackage{mathptmx}  % Times, aceptado por APA 7
\usepackage[spanish]{babel}
\usepackage[margin=2.54cm]{geometry}
\usepackage{setspace}
\usepackage{csquotes}          % antes de biblatex
\usepackage{graphicx}
\usepackage{booktabs}
\usepackage{array}
\usepackage{caption}
\usepackage{enumitem}
\usepackage{xcolor}
\usepackage{fancyhdr}
\usepackage[hidelinks]{hyperref}

% ---------- Bibliografía: APA 7 con biblatex-apa ----------
\usepackage[style=apa,backend=biber,language=spanish]{biblatex}
\DeclareLanguageMapping{spanish}{spanish-apa}
\addbibresource{referencias.bib}

% ---------- Formato APA 7 (documento estudiantil) ----------
\doublespacing                 % interlineado doble
\setlength{\parindent}{1.27cm} % sangría de 0.5 in
\setlength{\parskip}{0pt}
\setlist{nosep}
\captionsetup{font=small,labelfont=bf,justification=raggedright,singlelinecheck=false}

\setlength{\headheight}{14.5pt}
\pagestyle{fancy}
\fancyhf{}
\fancyhead[R]{\thepage}
\renewcommand{\headrulewidth}{0pt}

% Encabezados de nivel 1-3 al estilo APA 7
\usepackage{titlesec}
\titleformat{\section}{\normalfont\normalsize\bfseries\centering}{}{0pt}{}
\titleformat{\subsection}{\normalfont\normalsize\bfseries}{}{0pt}{}
\titleformat{\subsubsection}{\normalfont\normalsize\bfseries\itshape}{}{0pt}{}
\titlespacing*{\section}{0pt}{0pt}{0pt}
\titlespacing*{\subsection}{0pt}{0pt}{0pt}
\titlespacing*{\subsubsection}{0pt}{0pt}{0pt}

% Marcador de texto pendiente
\newcommand{\relleno}[1]{\textcolor{red}{[#1]}}

% Título y metadatos (usados en portada.tex)
\newcommand{\titulo}{Física computacional y gestión del riesgo de desastres en el Perú y Latinoamérica: aplicaciones, brechas de acceso y limitaciones}
\newcommand{\autor}{\relleno{Nombre y apellidos}}
\newcommand{\correo}{ccapatinta@unsa.edu.pe}

\begin{document}

% ---------- 1) Título y 2) integrantes con e-mails ----------
% Portada estilo APA 7 (estudiantes) adaptada a UNSA
\begin{titlepage}
  \singlespacing
  \centering
  \vspace*{3\baselineskip}
  {\bfseries \titulo \par}
  \vspace{2\baselineskip}
  \autor\par
  \correo\par
  \vspace{1\baselineskip}
  Escuela Profesional de Ingeniería de Sistemas\par
  Facultad de Ingeniería de Producción y Servicios\par
  Universidad Nacional de San Agustín de Arequipa\par
  \vspace{1\baselineskip}
  Organización y Métodos\par
  Mg. William Bornás Rios\par
  \vspace{1\baselineskip}
  \relleno{Fecha de entrega}\par
  Arequipa, Perú
\end{titlepage}

% ---------- 3) Resumen (mín. 250 palabras) y 4) Palabras clave ----------
\newpage
\begin{center}\textbf{Resumen}\end{center}
\noindent\relleno{Resumen de mínimo 250 palabras: contexto, objetivo, método (bases, query, criterios, número de registros), hallazgos por RQ y conclusión principal. Sin citas.}

\vspace{0.5\baselineskip}
\noindent\textbf{\textit{Palabras clave:}} \textit{\relleno{física computacional, simulación numérica, gestión del riesgo de desastres, responsabilidad social, Perú, Latinoamérica}}

\vspace{1\baselineskip}
\begin{center}\textbf{Abstract}\end{center}
\noindent\relleno{Opcional: versión en inglés del resumen.}

\vspace{0.5\baselineskip}
\noindent\textbf{\textit{Keywords:}} \textit{\relleno{computational physics, numerical simulation, disaster risk management, social responsibility, Peru, Latin America}}

% ---------- 5) Introducción (aprox. 8 párrafos medianos) ----------
\newpage
\begin{center}\textbf{\titulo}\end{center}
\section{Introducción}
\relleno{P1. Contexto: desastres naturales y de origen humano en el Perú y Latinoamérica.}

\relleno{P2. Física computacional: qué es y por qué importa para la gestión del riesgo.}

\relleno{P3. Responsabilidad social: impacto de la tecnología y bien común (deber social).}

\relleno{P4. Brechas de acceso, capacitación y comunicación en poblaciones vulnerables.}

\relleno{P5. Conexión con Organización y Métodos: cómo las organizaciones públicas y privadas coordinan procesos, roles y decisiones para gestionar el riesgo.}

\relleno{P6. Justificación y vacío en la literatura.}

\relleno{P7. Objetivo y objeto de estudio.}

\relleno{P8. Preguntas de investigación (RQ1--RQ3) y estructura del artículo.}

% ---------- 6) Background (máx. 2 páginas, imagen obligatoria) ----------
\section{Marco conceptual (\textit{background})}
\relleno{Conceptos básicos: física computacional, simulación numérica, modelación hidrodinámica, HPC, aprendizaje automático informado por la física, gestión del riesgo de desastres, responsabilidad social. Máximo 2 páginas.}

\begin{figure}[ht]
  \centering
  % \includegraphics[width=0.8\textwidth]{figuras/background.png}
  \fbox{\parbox{0.8\textwidth}{\centering\relleno{Imagen obligatoria del background}}}
  \caption{\relleno{Título de la figura}}
  \label{fig:background}
  \par\smallskip\noindent\small\textit{Nota.} \relleno{Fuente o adaptación de la figura, con cita.}
\end{figure}

% ---------- 7) Revisión sistemática ----------
\section{Revisión sistemática}

\subsection{Preguntas de investigación}
\begin{itemize}
  \item \textbf{RQ1 (aplicaciones):} ¿Qué técnicas de física computacional (simulación, modelado, HPC, aprendizaje automático informado por la física) se han aplicado a la prevención de desastres naturales en el Perú y Latinoamérica?
  \item \textbf{RQ2 (acceso y alcance):} ¿Qué barreras de acceso, capacitación y comunicación enfrentan las poblaciones vulnerables para beneficiarse de estas herramientas?
  \item \textbf{RQ3 (actores y proyección):} ¿Qué instituciones, organizaciones y programas formativos impulsan estas aplicaciones, y qué perspectivas de desarrollo se reportan?
\end{itemize}

\subsection{Estrategia de búsqueda}
\relleno{Bases: Scopus y SciELO. Fecha de búsqueda: 22/10/2026. Rango: 2021--2026. Idiomas: español e inglés. Tipos: artículos y revisiones. Query exacto de cada base.}

\begin{table}[ht]
  \centering
  \caption{Cadenas de búsqueda utilizadas}
  \label{tab:queries}
  \begin{tabular}{@{}p{2.5cm}p{9cm}r@{}}
    \toprule
    Base & Cadena de búsqueda & Resultados \\
    \midrule
    Scopus & \relleno{query} & 86 \\
    SciELO & \relleno{query unificado} & \relleno{n} \\
    \bottomrule
  \end{tabular}
  \par\smallskip\noindent\small\textit{Nota.} \relleno{Aclaraciones sobre filtros aplicados.}
\end{table}

\subsection{Criterios de inclusión y exclusión}
\textbf{Inclusión:}
\begin{itemize}
  \item Revisados por pares, 2021--2026, en español o inglés.
  \item Aplican o discuten física computacional en gestión de riesgo de desastres.
  \item Contexto de Perú o Latinoamérica.
  \item Aportan evidencia para al menos una RQ.
\end{itemize}
\textbf{Exclusión:}
\begin{itemize}
  \item Duplicados.
  \item Sin relación con desastres.
  \item Sin texto completo accesible.
  \item Solo teóricos, sin aplicación o caso concreto.
\end{itemize}

\subsection{Selección de estudios}
\relleno{Diagrama de flujo tipo PRISMA y explicación: identificados $\rightarrow$ duplicados eliminados $\rightarrow$ cribado por título y resumen $\rightarrow$ texto completo $\rightarrow$ incluidos.}

\begin{table}[ht]
  \centering
  \caption{Flujo de selección de estudios}
  \label{tab:flujo}
  \begin{tabular}{@{}lrrr@{}}
    \toprule
    Etapa & Scopus & SciELO & Total \\
    \midrule
    Registros identificados & 86 & \relleno{n} & \relleno{n} \\
    Tras eliminar duplicados & & & \relleno{n} \\
    Tras cribado por título y resumen & & & \relleno{n} \\
    Evaluados a texto completo & & & \relleno{n} \\
    Incluidos en el análisis final & & & \relleno{n} \\
    \bottomrule
  \end{tabular}
\end{table}

\subsection{Análisis por pregunta de investigación}
\subsubsection{RQ1: aplicaciones}
\relleno{Análisis con evidencia (citas \textbackslash parencite).}
\subsubsection{RQ2: acceso y alcance}
\relleno{Análisis con evidencia.}
\subsubsection{RQ3: actores y proyección}
\relleno{Análisis con evidencia.}

% ---------- 8) Limitaciones y consideraciones ----------
\section{Limitaciones y consideraciones}
\relleno{Pocos registros con Perú como caso de estudio en SciELO; exclusión de portugués; solo literatura revisada por pares; rango temporal; sesgo de las palabras clave.}

% ---------- 9) Conclusiones ----------
\section{Conclusiones}
\relleno{Respuesta sintética a RQ1--RQ3, aportes y recomendaciones (por ejemplo, una app o herramienta para quienes no reciben alertas nacionales).}

% ---------- 10) Referencias (mín. 18) ----------
\newpage
\begin{center}\textbf{Referencias}\end{center}
\printbibliography[heading=none]

\end{document}